\chapter{Appendix A: Glossary}

This glossary is intentionally pragmatic.
It captures the terms used throughout the book in the sense they are used in this repository.

\section{VNode / VDOM}
A \textbf{VNode} (virtual node) is an in-memory object that describes what the UI should be.
A \textbf{VDOM tree} is a nested structure of VNodes.

Why it matters: a VDOM lets you describe UI declaratively, then reconcile the description with the real DOM.

\section{Mount}
\textbf{Mounting} is turning a VNode into real DOM nodes and inserting them into the document.

\section{Patch}
\textbf{Patching} is updating an existing DOM subtree to match a new VNode tree.
It\textquotesingle s the \emph{heart} of a framework: correctness first, then performance.

\section{Unmount / Destroy}
\textbf{Unmounting} (or destroying) removes DOM nodes and cleans up resources (event listeners, subscriptions, effects).

\section{Keyed children}
When a list of children has \textbf{keys}, the renderer can preserve identity across reorders.
Keys are about stable meaning, not indexes.

\section{Scheduler}
A \textbf{scheduler} decides \emph{when} work runs.
It is the boundary between synchronous user code and asynchronous rendering.

\section{Signals / Effects}
\textbf{Signals} represent reactive state.
\textbf{Effects} are computations that re-run when signals they read change.

\section{HRBR}
In this book, \textbf{HRBR} refers to the compiler-assisted rendering approach used by Relax.js.
The compiler tries to extract a stable HTML template and a set of dynamic \textbf{slots}.
At runtime, a mounted block updates only the slots.
When the compiler can\textquotesingle t statically guarantee structure, it emits a safe fallback.
